\chapter{Complexity — Directed Walk}
%%% generated by the blueprint pipeline from TCSlib.Complexity.Randomized.DirectedWalk
\section{Overview}

Declarations of the Lean module \texttt{TCSlib.Complexity.Randomized.DirectedWalk}.

\section{Declarations}

\begin{definition}[Outgoing neighbors]
\lean{Randomized.outNbrs}
\statementsource{Vad12}{
  pdf-bf2ce5422088-p038-b004,
  pdf-bf2ce5422088-p039-b002
}
\label{Randomized.outNbrs}
\leanok
In a finite directed graph, the outgoing neighbors of a vertex $v$ are the vertices $u$
such that there is a directed edge from $v$ to $u$.
\end{definition}

\begin{definition}[Average over outgoing neighbors]
\lean{Randomized.avgOut}
\statementsource{Vad12}{
  pdf-bf2ce5422088-p039-b002,
  pdf-bf2ce5422088-p039-b003
}
\label{Randomized.avgOut}
\leanok
\uses{Randomized.outNbrs}
Let $G$ be a finite directed graph, $f$ a rational-valued function on its vertices, and
$v$ a vertex. Define the outgoing average of $f$ at $v$ as the arithmetic mean of $f$
over the outgoing neighbors of $v$ when there is at least one; when there are none,
define it to be $f(v)$.
\end{definition}

\begin{definition}[Finite-step hitting values]
\lean{Randomized.hitProb}
\statementsource{Vad12}{pdf-bf2ce5422088-p038-b006}
\label{Randomized.hitProb}
\leanok
\uses{Randomized.avgOut}
Let $G$ be a finite directed graph and let $t$ be a vertex. Define $h_k(s)\in\mathbb{Q}$
for each $k\in\mathbb{N}$ and vertex $s$ recursively by $h_0(s)=1$ if $s=t$ and
$h_0(s)=0$ otherwise; for $k\geq 0$, set $h_{k+1}(s)=1$ if $s=t$, and otherwise let it
be the average of $h_k$ over the outgoing neighbors of $s$, or $h_k(s)$ when $s$ has no
outgoing neighbors.
\end{definition}

\begin{definition}[Truncated expected hitting time]
\lean{Randomized.hitTime}
\statementsource{Vad12}{pdf-bf2ce5422088-p038-b006}
\label{Randomized.hitTime}
\leanok
\uses{Randomized.hitProb}
Let $G$ be a finite directed graph, with target vertex $t$, starting vertex $s$, and
natural number $N$. Define the truncated hitting time as the sum, over $0 \le j < N$, of
one minus the probability of having hit $t$ within $j$ steps when starting at $s$.
\end{definition}

\begin{definition}[Directed chain graph]
\lean{Randomized.chainDigraph}
\statementsource{Vad12}{
  pdf-bf2ce5422088-p047-b004,
  pdf-bf2ce5422088-p047-b005
}
\label{Randomized.chainDigraph}
\leanok
For a natural number $n$, define a directed graph on the vertices $0,1,\ldots,n$. There
is an edge from $i$ to $j$ exactly when $j=0$ or $j=i+1$.
\end{definition}

\begin{theorem}[Reachability in a Chain Digraph]
\lean{Randomized.chainDigraph_reachable}
\label{Randomized.chainDigraph_reachable}
\leanok
\uses{Randomized.chainDigraph}
\difficulty{4}
For every natural number $n$, consider the directed graph on $\{0,1,\ldots,n\}$ in which
there is an edge from $i$ to $j$ exactly when $j=0$ or $j=i+1$. There is a directed
path, allowing a path of length zero, from $0$ to $n$.
\end{theorem}

\begin{theorem}[Average over chain out-neighbours]
\lean{Randomized.avgOut_chainDigraph}
\label{Randomized.avgOut_chainDigraph}
\leanok
\uses{Randomized.avgOut, Randomized.chainDigraph, Randomized.outNbrs}
\difficulty{4}
Let $n$ be a natural number, let $f:	ext{Fin}(n+1)	obr$, and let $i	ext{ in Fin}(n+1)$
satisfy $i<n$. In the digraph on $	ext{Fin}(n+1)$ where $j$ is an out-neighbour of $i$
exactly when $j=0$ or $j=i+1$, the average of $f$ over the out-neighbours of $i$ is
$(f(0)+f(i+1))/2$.
\end{theorem}

\begin{theorem}[Pointwise hitting-probability bound on a chain]
\lean{Randomized.hitProb_chainDigraph_le_aux}
\label{Randomized.hitProb_chainDigraph_le_aux}
\leanok
\uses{Randomized.avgOut_chainDigraph, Randomized.chainDigraph, Randomized.hitProb}
\difficulty{6}
For all natural numbers $n$ and $k$ and every vertex $i\in\text{Fin}(n+1)$, consider the
digraph on $text{Fin}(n+1)$ with an edge from $u$ to $v$ when $v=0$ or $v=u+1$. The
$k$-step probability of hitting the target vertex $n$ when starting at $i$, multiplied
by $2^n$, is at most $2^i+k$.
\end{theorem}

\begin{theorem}[Hitting probability on a chain digraph]
\lean{Randomized.hitProb_chainDigraph_le}
\label{Randomized.hitProb_chainDigraph_le}
\leanok
\uses{Randomized.chainDigraph, Randomized.hitProb, Randomized.hitProb_chainDigraph_le_aux}
\difficulty{3}
For all natural numbers $n$ and $T$, consider the digraph on $	ext{Fin}(n+1)$ in which
$i$ has an edge to every vertex $j$ that is $0$ or satisfies $j=i+1$. Starting at vertex
$0$, the probability of hitting the target vertex $n$ within the recursively defined
$T$-step process is at most $(T+1)/2^n$.
\end{theorem}

\begin{theorem}[Sum of consecutive integers]
\lean{Randomized.sum_range_succ_cast}
\label{Randomized.sum_range_succ_cast}
\leanok
\difficulty{4}
For every natural number $N$, the sum of $j+1$ over $0\le j<N$, interpreted in the
rationals, equals $N(N+1)/2$.
\end{theorem}

\begin{theorem}[Lower bound on chain hitting time]
\lean{Randomized.hitTime_chainDigraph_ge}
\label{Randomized.hitTime_chainDigraph_ge}
\leanok
\uses{Randomized.chainDigraph, Randomized.hitProb_chainDigraph_le, Randomized.hitTime, Randomized.sum_range_succ_cast}
\difficulty{5}
For every natural number $n$ and rational number $B$, if the expected truncated hitting
time of the last vertex in the chain digraph on $	ext{Fin}(n+1)$, starting from vertex
$0$, is at most $B$ for every truncation length $N$, then $(2^n-1)/2\le B$.
\end{theorem}
